% The preface, in the author's voice; \input by build.mjs.

Suppose every pixel of a camera measured, instead of the light that falls on it, the tide of gravity at its place:
the small difference in pull between one side of the pixel and the other. What would such a camera see? It would
see through walls, because nothing screens gravity. It would weigh light, because all energy gravitates. It would
see the Sun at midnight, shining up through the dark Earth. And it would be blind in ways that a camera of light is
not: a shell and a ball of the same mass would look the same to it, and its pictures would blur in a way that no lens
can correct.

This book answers the question in full: what such an instrument sees, how sharply, with how much noise, and what no
instrument of its kind can ever see. It is written for a reader who knows calculus, vectors and matrices and a
little probability, and who is willing to work. It builds from there all the physics the answer
needs: Newton's gravity and its tides, harmonic functions and multipoles, noise and estimation, oscillators and
heat, waves and optics, special relativity, curvature and Einstein's equation, gravitational waves, quantum
mechanics as far as the limits it sets on measurement, black holes, and the sky. Each of these is developed for
its own sake, as a textbook would develop it, and each is then used. Every result the book states is derived in
it, or, where it comes from the literature, stated with its source.

The camera's own results come at the end of each part, in a chapter called after what the camera sees there, and
they gather in Part~VI. Many of them are new, and each is derived where it is stated.

The book is long because the subject is wide, not because any one part is difficult. It can be read in order, which
is how it is meant to be read. A reader who knows some of the physics can skip forward; the sections say what they
use, and the cross-references say where it was derived. The exercises at the end of each chapter come in tiers:
practice, which checks the section; stretch, which takes it further; and problems drawn from the camera itself.
Selected solutions are at the back.

The book has a web edition, in which the figures move. Most of the figures in these pages are frames of
interactive drawings there, and a margin note says where one exists.

Every number in the book was computed by a reference implementation, a small program written for the purpose and
tested against independent calculations; Appendix~B describes it and says what each test establishes. The figures
were drawn by code from the same numbers.

This book was written with Claude (Anthropic), used throughout as a research and writing
assistant. Under my direction it drafted and revised the text, carried out and checked derivations, searched the
literature, and wrote the reference implementation, its tests and the code that draws the figures. The question,
the choice of what the book says and the judgement of its results are mine, and I take responsibility for every
statement in it. The cover art and the pictured objects in some figures were made with an image model; every mark
that carries a number is drawn by code from computed results.

\medskip
\noindent\hfill Mostafa Mohsenvand\\
\hspace*{\fill}Cambridge, Massachusetts, 2026
